% !PRINTED 348
\weberpage{348}

%<<CONT
\gesperrt{is determined.} There must therefore be four rational integers
$\alpha$, $\beta$, $\gamma$, $\delta$ which satisfy the condition
\weq{2}{\begin{aligned}\eta\, \mathfrak{a}' &= \mathfrak{a}\,(\alpha - \gamma\,\omega),\\ \eta\, \mathfrak{a}'\,\omega' &= \mathfrak{a}\,(-\beta + \delta\,\omega)\end{aligned}}
If we denote for the moment by $\eta_1$, $\omega_1$, $\omega_1'$ the numbers
conjugate to $\eta$, $\omega$, $\omega'$, then accordingly
\[
\mathfrak{a}'^2\, \eta\, \eta_1
\begin{vmatrix} 1, & \omega' \\ 1, & \omega_1' \end{vmatrix}
= \mathfrak{a}^2
\begin{vmatrix} \alpha, & -\gamma \\ -\beta, & \delta \end{vmatrix}
\begin{vmatrix} 1, & \omega \\ 1, & \omega_1 \end{vmatrix},
\]
hence:
\[
\mathfrak{a}'^2\, \eta\, \eta_1\,(\omega' - \omega_1') = \mathfrak{a}^2\,(\alpha\delta - \beta\gamma)\,(\omega - \omega_1),
\]
and since, by \S\,91, (12):
\[
\mathfrak{a}\,(\omega - \omega_1) = \mathfrak{a}'\,(\omega' - \omega_1') = -\sqrt{\Delta}
\]
holds, there follows from this, if
\[
\varepsilon = \alpha\delta - \beta\gamma
\]
is put,
\weq{3}{\mathfrak{a}\,\varepsilon = \mathfrak{a}'\,\eta\,\eta_1;}
if therefore, as we have already assumed, $\mathfrak{a}$ and $\mathfrak{a}'$ are positive
and (according to the sharpened concept of equivalence) $\eta\eta_1$ is
positive, then $\varepsilon$ is positive as well.

But if we interchange $\mathfrak{a}$ with $\mathfrak{a}'$, then $\eta$ passes into $1 : \eta$, and
$\varepsilon$ may pass into $\varepsilon'$. Accordingly there follows also
\weq{4}{\mathfrak{a}'\,\varepsilon' = \frac{\mathfrak{a}}{\eta\,\eta_1},}
and consequently $\varepsilon \varepsilon' = 1$. But since $\varepsilon$ and $\varepsilon'$ are positive integral
rational numbers, it follows from this that $\varepsilon = \varepsilon' = 1$. Hence by (2)
\weq{5}{\omega' = \frac{-\beta - \delta\,\omega}{\alpha - \gamma\,\omega}, \qquad \omega = \frac{\alpha\,\omega' + \beta}{\gamma\,\omega' + \delta},}
\[
\alpha\delta - \beta\gamma = 1.
\]

Now $\omega$ and $\omega'$ are the roots of the two forms
\[
(a,\, b,\, c), \quad (a',\, b',\, c') \qquad [\S\,91,\ (13)],
\]
and the equivalence of these forms is expressed by (5).
We therefore have, in agreement with the general theory
(\S\,88), the theorem:

5. \gesperrt{The quadratic forms belonging to equivalent ideals
are equivalent.}

From this one recognises the significance of the sharpened concept of
equivalence: for if one wished to admit negative values of
$N(\eta)$ as well, then $\varepsilon = -1$ would also be possible, and the
corresponding forms could be \gesperrt{improperly} equivalent.
